\section{Integral local types over the original coefficient ring}
\label{sec:integral-type-v43}

\status{This chapter supplies the original integral
\texttt{TypeForAlt} input to OTW13.  The construction keeps the
given coefficient ring, residual characters, Galois quotient,
local groups and complex coefficient map.  Its only new substantial
published input is the crystalline inertial-type theorem, at the
precise scope stated in ITA6.  We prove the descent to the original
ring and the universal local-correspondence and Teichm\"uller
quantifiers.  Classical local Langlands and the original
nonemptiness theorem \texttt{RecTExists} remain foundational inputs;
the latter is used only to prevent a universal quantifier over
correspondences from being vacuous.  No global automorphy,
nonvanishing, Hida endpoint, or RH/GRH conclusion is asserted here.}

Fix the original prime $\ell$, CM field $L$, algebraic closure
$\Omega$, setting $s:\texttt{HidaSetting}$, integer $n\geq1$,
level datum $V:s.\texttt{Lev}$, auxiliary datum $Q$, coefficient
datum $c$, residual datum $r$, and field isomorphism
\[
 \iota:\overline{\mathbb Q}_\ell\xrightarrow{\sim}\mathbb C.
\]
Write $O=c.O$, $K=\operatorname{Frac}(O)$,
$k=O/\mathfrak m_O$, and $q_O=|k|$.  The given coefficient
map is
\[
 j:O\lhook\joinrel\longrightarrow\overline{\mathbb Q}_\ell
              \xrightarrow{\iota}\mathbb C.
 \tag{ITA.1}
\]
We also write $j$ for its extension to $K$; neither $j$ nor
$\iota$ is assumed continuous for the complex topology.
The original residual characters are
$\bar\chi_i=r.\texttt{chib}_i:G_0\to k^\times$,
$0\leq i<n$.  Their kernels are open.  Their pairwise
distinctness, polarization and the other fields of $r$ will not
be needed for this construction.

The finite set $S=s.\texttt{St}$ is the source's set of type
places.  For $u\in S$, put
\[
 F_u=L_{s.w(u)},\quad G_u=\operatorname{GL}_n(F_u),\quad
 K_u=\operatorname{GL}_n(\mathcal O_{F_u}),\quad
 D_u=\operatorname{Dec}_L(s.A(u)),\quad
 I_u=\operatorname{Ine}_L(s.A(u)).
 \tag{ITA.2}
\]
These are the actual groups at the specified valuations.
The original map $V.\mathrm{pr}:G_L\to G_0$ restricts to
each $D_u$ and $I_u$.  Write
$G_S=\prod_{u\in S}G_u$ and $K_S=\prod_{u\in S}K_u$;
these identify with the source's \texttt{GV} and \texttt{KT}
by the coordinate inclusions.  The residue characteristic of
$F_u$ can differ from $\ell$; in this setting it does differ.
We will impose no bound comparing either prime to $n$.

Here is the exact conclusion.  There are an integer $m$ and one
homomorphism
\[
                    \tau_0:K_S\longrightarrow\operatorname{GL}_m(O)
 \tag{ITA.3}
\]
such that its complexification $\tau_{0,j}$ has open kernel and
the following equivalence holds.  It holds for every family
$t_i:G_0\to O^\times$ satisfying
\[
 \overline{t_i(g)}=\bar\chi_i(g),\qquad
                       t_i(g)^{q_O-1}=1,                    \tag{ITA.4}
\]
every nonzero complex vector space $M$, every representation
$\rho:G_S\to\operatorname{Aut}_{\mathbb C}(M)$, and every
family of smooth irreducible representations $\pi_u$ for which
$M$, restricted to $G_u$, is generated by the images of all
$G_u$-equivariant maps from $\pi_u$.  The last condition is the
literal \texttt{IsoSum} condition; no finite length or
admissibility of $M$ is added.  Define
\[
 \chi^D_{u,i}(g)=j\bigl(t_i(V.\mathrm{pr}(g))\bigr),\qquad
                              \chi_{u,i}=\chi^D_{u,i}|_{I_u}.
 \tag{ITA.5}
\]
The asserted equivalence is
\begin{align}
 &\exists (F_a)_{0\leq a<m}\in M^m\setminus\{0\},\quad
   \sum_{b=0}^{m-1}j(\tau_0(k)_{ab})\,\rho(k)F_b=F_a
                 \quad(k\in K_S,\ 0\leq a<m)
                                                        \tag{ITA.6}\\
 &\hspace{25mm}\Longleftrightarrow\quad
   \forall u\in S\ \forall\mathrm{rc},\quad
   \texttt{IsRecT}(s.hw(u),\mathrm{rc})\Longrightarrow
   \texttt{IsTypeWD}\bigl(\chi_u,
                       \mathrm{rc}(\operatorname{algNorm}\pi_u)\bigr).
                                                        \notag
\end{align}
Here
$\operatorname{algNorm}\pi=\pi\otimes|\det|^{(1-n)/2}$.
The source predicate \texttt{IsTypeWD} asks for a representative
of its Weil--Deligne isomorphism class on $\mathbb C^n$ with
$N=0$ and inertia matrices exactly
$\operatorname{diag}(\chi_0,\ldots,\chi_{n-1})$.  Repeated
inertia characters are allowed.  Formula (ITA.6), including the
order of all its quantifiers, is \texttt{TypeForAlt}.

\subsection{The actual coefficient ring and every Teichm\"uller lift}

% BEGIN ITA1
\paragraph{ITA1. Completeness and the algebraic coefficient closure.}
\label{claim:ITA1}
The ring $O$ is the complete valuation ring of the finite extension
$K/\mathbb Q_\ell$, its residue field is finite, and
$\mathbb C$ is an algebraic closure of $j(K)$.
\begin{proof}
GA1 identifies the actual subring $O\subset\overline{\mathbb Q}_\ell$
with the elements of $K$ of norm at most one, and GA2 proves
that $K$ is finite over $\mathbb Q_\ell$.  The equivalence of
the coordinate and induced norms in GA2 proves that $K$ is
complete.  Its closed unit ball $O$ is therefore complete.
The norm topology and the maximal-ideal topology agree by
GA1's uniformizer formula.  Since $O$ is a finite
$\mathbb Z_\ell$-module, $O/\ell O$ is finite; $k$ is its
quotient because $\ell\in\mathfrak m_O$.  Finally every
element of $\overline{\mathbb Q}_\ell$ is algebraic over
$\mathbb Q_\ell\subset K$, and $\iota$ is onto the
algebraically closed field $\mathbb C$.  Thus every complex
number is algebraic over this particular embedded copy $j(K)$.
This is an algebraic statement and supplies no complex
continuity assertion.
\end{proof}
% END ITA1

% BEGIN ITA2
\paragraph{ITA2. One and only one family satisfies the original lift predicate.}
\label{claim:ITA2}
There is a unique multiplicative lift $[\ ]:k^\times\to O^\times$
with $[a]^{q_O-1}=1$ and residue $a$.  Hence (ITA.4) has the
unique solution $t_i=[\ ]\circ\bar\chi_i$, and
\[
                         \ker t_i=\ker\bar\chi_i.       \tag{ITA.7}
\]
\begin{proof}
For $a\in k^\times$, choose any unit lift $a_0\in O$.
The polynomial $f(X)=X^{q_O-1}-1$ vanishes at $a_0$
modulo $\mathfrak m_O$, and
$f'(a_0)=(q_O-1)a_0^{q_O-2}$ is a unit: the residue of
$q_O-1$ is $-1$.  Newton iteration
$a_{r+1}=a_r-f(a_r)/f'(a_r)$ stays in the same residue
class, keeps the derivative invertible, and at least doubles
the valuation of $f(a_r)$.  Its successive differences tend
to zero geometrically in the maximal-ideal filtration.
Completeness from ITA1 gives a root $[a]$ in that residue
class.  This also proves existence directly at $\ell=3$.

For uniqueness let $x,y\in O^\times$ have that residue and
$(q_O-1)$st power one.  Put $z=xy^{-1}$, so $\bar z=1$.
Then
\[
          (z-1)(1+z+\cdots+z^{q_O-2})=z^{q_O-1}-1=0.
\]
The second factor is a unit, since its residue is $-1$.
Consequently $z=1$.  Products of the chosen lifts satisfy
the same equation and have the product residue, proving
$[ab]=[a][b]$; also $[1]=1$.  Applying uniqueness to each
value proves uniqueness of the entire family $t_i$ among
all homomorphisms in the source predicate.  A value $t_i(g)$
equals one exactly when its residue equals one, which gives
(ITA.7).  All kernels are therefore open, and all values
belong to the finite group of $(q_O-1)$st roots of unity.
\end{proof}
% END ITA2

% BEGIN ITA3
\paragraph{ITA3. The precise local inertia types and their field of values.}
\label{claim:ITA3}
For every $u$, the characters (ITA.5) on the original $D_u$
have finite image and open kernel.  They have values in
$j(O^\times)$, and
$\vartheta_u=\bigoplus_{i=0}^{n-1}\chi_{u,i}$ is an
inertial type for $F_u$.
\begin{proof}
The map $V.\mathrm{pr}$ is the original continuous restriction
and quotient map.  By (ITA.7), the kernel of its composite
with $t_i$ is the inverse image of the open subgroup
$\ker\bar\chi_i$.  The embedding $j$ is injective, so
pushing coefficients through $j$ does not change this
kernel.  The image is finite by ITA2, proving all the stated
topological assertions without any continuity of $j$.
The value assertion is the displayed formula (ITA.5) itself.

Complete at the fixed valuation $s.A(u)$.  Restriction and
completion of finite extensions identify $D_u,I_u$ and its
Weil subgroup with $G_{F_u},I_{F_u}$ and $W_{F_u}$, carrying
the specified subgroup inclusions to the usual ones.  This
is the same completion identification used in KX3 and is
valid at every finite place.  The direct sum of the
$\chi^D_{u,i}|_{W_{F_u}}$ extends $\vartheta_u$ to the
Weil group and has finite image and open kernel there.
Thus it meets the published definition of inertial type.
Neither distinctness of the $\chi_{u,i}$ nor a tame-image
condition is required.
\end{proof}
% END ITA3

\subsection{The local-correspondence dictionary, including its quantifiers}

For one place write $F=F_u$, $D=D_u$, $I=I_u$, $W=W_{F_u}$,
and $\chi_i^D=\chi^D_{u,i}$.  Let $\mathrm{rec}$ denote
the usual complex local Langlands correspondence with geometric
reciprocity.  Retain its rank-one formula and its compatibility
with ordered Langlands quotients.  For an ordered list of smooth
characters $\mu_0,\ldots,\mu_{n-1}$ satisfying
\[
                  i<j\ \Longrightarrow\ \mu_j\ne\mu_i|\ |,
                                                               \tag{ITA.8}
\]
normalized induction has a unique irreducible quotient; its
parameter is the sum of the corresponding Weil characters with
$N=0$.  These are the rank-one part of the usual Langlands
classification, not a claim that every principal series is
irreducible.  They are also precisely the \texttt{one},
\texttt{langlands} and singleton-segment laws used below.
One primary reference for these conventions and the quotient
statement is Caraiani--Emerton--Gee--Geraghty--Pa\v{s}k\={u}nas--Shin,
\emph{Patching and the $p$-adic local Langlands correspondence},
\S\,1.8, \S\,3.1 and Proposition 3.9, with its proof via
Bernstein--Zelevinsky classification.
\footnote{Author's full paper:
\url{https://www.ma.imperial.ac.uk/~tsg/Index_files/Gpatching.pdf}.
The local statement is independent of the global patching
restriction in that paper.  For the unrestricted inertial-type
theorem itself we use the separate statement in ITA6.}

The source has a larger class of maps, called \texttt{IsRecT}.
We use exactly the following fields of such a map: a map
$a:W\to F^\times$ satisfying \texttt{IsLocalArtin}; the
rank-one formula; the ordered-segment formula; compatibility
with real unramified twists; and injectivity against a
\texttt{Good} representation.  Representations with a
cuspidal support consisting of rank-one characters are
\texttt{Good} by the first disjunct of its definition.
The source's theorem \texttt{RecTExists} supplies at least
one such map at the given place.  Its exact foundational
endpoint is
\path{Foundations/LocalLanglands/LocalCentralUnits/RecTExists.lean},
and the caller is \texttt{WinConv.exRec'\_alt}.  This
nonemptiness is retained as part of the original local
Langlands foundation.  Our argument does not replace
\texttt{IsRecT} by a preferred map or assume its maps agree
on arbitrary representations.

% BEGIN ITA4
\paragraph{ITA4. Every allowed local Artin map at a type place is the geometric one.}
\label{claim:ITA4}
Every map $a:W\to F^\times$ satisfying the source's
\texttt{IsLocalArtin} equals the inverse of geometric local
reciprocity on the Weil abelianization.
\begin{proof}
The $u\in s.\mathrm{St}$ branch of $s.\mathrm{split}$
says that the original place below $w=s.w(u)$ splits in
$L/L^+$.  The same decomposition-group argument as in
CIP3 therefore gives $w\ne\bar w$.  In particular CIP3's
local-unit injectivity applies to this place, although
it lies away from $\ell$.

Put $a_0=(\operatorname{rec}^{\mathrm{geo}}_F)^{-1}$ on
$W^{\mathrm{ab}}$, using the actual completion in ITA3.
Local class field theory gives its surjectivity, continuity
on inertia, and
$\operatorname{ord}_F(a_0(\sigma))=-\deg(\sigma)$,
where degree one denotes arithmetic Frobenius.  A smooth
Weil character factors through this abelianization and
therefore through $a_0$, giving the source's \texttt{factor}
field.  The resulting character of $F^\times$ has open
kernel, because the reciprocity quotient is open.
Local--global compatibility gives
\[
       \alpha_L(j_w a_0(\sigma))=\operatorname{toAb}(\sigma),
                                                               \tag{ITA.9}
\]
where $\alpha_L$ is the inverse-arithmetic global map of
CIP1.  CIP2 identifies every source global Artin map with
$\alpha_L$, so this proves the universally quantified
\texttt{glob} field as well.

For any other source map $a$, the quotient
$a(\sigma)/a_0(\sigma)$ is a unit by \texttt{ord\_eq}.
Its local idele is killed by $\alpha_L$, by the
\texttt{glob} field and (ITA.9).  CIP3 forces that unit
to be one.  This holds for every $\sigma\in W$ and
proves uniqueness.  It also checks the reciprocity sign:
a geometric Frobenius maps to a uniformizer.  No
cyclotomic-prime hypothesis has entered this proof.
\end{proof}
% END ITA4

% BEGIN ITA5
\paragraph{ITA5. Exact equivalence with the every-\texttt{IsRecT} condition.}
\label{claim:ITA5}
For every smooth irreducible $\pi$ of $\operatorname{GL}_n(F)$,
the following three conditions are equivalent:
\begin{align*}
 P_\chi(\pi):\quad &N=0\text{ and }
     \mathrm{rec}(\operatorname{algNorm}\pi)|_I
                                  \simeq\bigoplus_i\chi_i;\\
 U_\chi(\pi):\quad &\texttt{IsTypeWD}
       (\chi,\mathrm{rc}(\operatorname{algNorm}\pi))
                     \text{ for every }\mathrm{rc}\text{ satisfying }
                     \texttt{IsRecT};\\
 E_\chi(\pi):\quad &\text{the same assertion for at least one such }
                     \mathrm{rc}.
\end{align*}
No global uniqueness theorem for the weaker source correspondence
is needed for this equivalence.
\begin{proof}
First consider a Frobenius-semisimple Weil--Deligne
representation $(R,N)$ with $N=0$ and
$R|_I\simeq\bigoplus_i\chi_i$, where every $\chi_i$
extends to the character $\chi_i^D$ of $D$.
Each occurring inertia character is fixed by conjugation
by $W$, since its extension is a character.  Consequently
every inertia-isotypic subspace is preserved by a
Frobenius lift $\Phi$.  On that subspace inertia acts
scalarly and $R(\Phi)$ is semisimple.  Choose an eigenbasis
there.  Since $W/I\simeq\mathbb Z$ is generated by
$\Phi$, the resulting basis splits the entire Weil
representation into smooth characters
\[
                            R=\bigoplus_{i=0}^{n-1}\xi_i.
                                                               \tag{ITA.10}
\]
Their inertia restrictions, with multiplicities, are exactly
the original $\chi_i$.  The inertia image is finite, so
every $\xi_i$ is smooth in the source sense.  Reordering
the basis to the original inertia list gives the literal
diagonal matrices demanded by \texttt{IsTypeWD}; this
also proves its equivalence with the stated $N=0$
and inertia-isomorphism conditions.

By local reciprocity, or the \texttt{factor} field,
$\xi_i=\mu_i\circ a_0$ for smooth $\mu_i:F^\times\to
\mathbb C^\times$.  Reorder them by nondecreasing
$|\mu_i(\varpi_F)|$.  Since $|\varpi_F|<1$, this order
satisfies (ITA.8).  Let $\Pi$ be the unique irreducible
quotient of their normalized induction.  It has rank-one
cuspidal support and is therefore \texttt{Good}.  Its
parameter under the usual correspondence is (ITA.10),
with zero monodromy.  More significantly, for \emph{every}
source map $\mathrm{rc}'$, ITA4 identifies its Artin map
with $a_0$, and its \texttt{one} and \texttt{langlands}
fields give the same direct sum (ITA.10).

If $E_\chi(\pi)$ holds, apply this construction to
$\mathrm{rc}(\operatorname{algNorm}\pi)$ for its witnessing
map.  That map sends $\Pi$ to the same parameter.  Its
\texttt{inj} field, whose second argument $\Pi$ is
\texttt{Good}, gives
$\operatorname{algNorm}\pi=\Pi$ as irreducible classes.
Every other source map and the usual correspondence
therefore give the required type.  Thus
$E_\chi\Rightarrow U_\chi$ and $E_\chi\Rightarrow P_\chi$.

If $P_\chi(\pi)$ holds, construct $\Pi$ from its usual
parameter.  The bijectivity of classical local Langlands
gives $\operatorname{algNorm}\pi=\Pi$; the preceding
rank-one calculation gives $U_\chi(\pi)$ and, using
\texttt{RecTExists}, $E_\chi(\pi)$.  Finally
$U_\chi\Rightarrow E_\chi$ uses exactly that same
nonemptiness theorem.  All three conditions are equivalent.

An unramified real twist of a Weil--Deligne representation
multiplies $R(w)$ by an unramified nonzero scalar and leaves
$N$ unchanged.  Hence $P_\chi(\pi)$ is also equivalent to
$N=0$ and the same inertia type for $\mathrm{rec}(\pi)$.
This checks explicitly the source's $(1-n)/2$ normalization
against the unnormalized form of the inertial-type theorem.
\end{proof}
% END ITA5

% BEGIN ITA6
\paragraph{ITA6. The published local type applies at the complete original scope.}
\label{claim:ITA6}
There is a nonzero finite-dimensional smooth irreducible complex
representation $A$ of $K_F=\operatorname{GL}_n(\mathcal O_F)$
such that for every smooth irreducible $\pi$ of
$\operatorname{GL}_n(F)$,
\[
             \operatorname{Hom}_{K_F}(A,\pi)\ne0
                         \quad\Longleftrightarrow\quad P_\chi(\pi).
                                                               \tag{ITA.11}
\]
Its kernel is open and it factors through a finite quotient of
$K_F$.
\begin{proof}[Published theorem and application]
Use the crystalline part of Matthew Emerton and Toby Gee,
\emph{Moduli stacks of \textup{\'etale} $(\varphi,\Gamma)$-modules
and the existence of crystalline lifts}, Annals of Mathematics
Studies 215, Theorem 8.2.1, with its inertial-type definition
in \S\,1.12.
\footnote{Author's text, printed p.~204 (PDF page 214),
and pp.~15--16 for conventions:
\url{https://www.math.uchicago.edu/~emerton/pdffiles/moduli.pdf}.}
For any finite extension $F/\mathbb Q_p$ and any inertial
type of rank $n$, it supplies a smooth irreducible
finite-dimensional representation $\sigma^{\mathrm{crys}}$
of $K_F$.  Its multiplicity in any irreducible smooth
representation is nonzero exactly for that inertia type
and zero monodromy.  There is no genericity hypothesis on
$\pi$, no restriction $p\nmid n$, and no exclusion of
$p=3$.  We need only existence and detection, not the
additional multiplicity bound.  This statement cites
CEGGPS, Theorem 3.7; using the book's unrestricted local
statement avoids importing that article's separate
global patching hypotheses.

Apply it to $\vartheta=\bigoplus_i\chi_i$, which is an
inertial type by ITA3.  The theorem is written over
$\overline{\mathbb Q}_p$; transport it to $\mathbb C$
through a field isomorphism.  This is an algebraic
transport of smooth representations, with no comparison
between complex and $p$-adic topologies.  Its usual
reciprocity convention is geometric, as in ITA4.
The possible unramified normalization is harmless by the
last paragraph of ITA5.  Thus its conclusion is exactly
(ITA.11), at the original $F$ and $n$.

Choose a finite basis of $A$.  The intersection of its
open stabilizers is contained in the kernel, so that
kernel is open.  It is normal, and an open subgroup
of compact $K_F$ has finite index: its cosets are an
open cover and compactness gives a finite subcover.
Thus $A$ factors through the finite quotient
$H=K_F/\ker A$.
\end{proof}
% END ITA6

% BEGIN ITA7
\paragraph{ITA7. Invariance under all algebraic automorphisms fixing the coefficient ring.}
\label{claim:ITA7}
For every $\gamma\in\operatorname{Aut}(\mathbb C/j(K))$
and every smooth irreducible $\pi$, one has
\[
                        P_\chi({}^{\gamma}\pi)
                              \Longleftrightarrow P_\chi(\pi).
                                                               \tag{ITA.12}
\]
Here scalar conjugation gives the original representation an
equivalent $\gamma$-semilinear model.  Arbitrary algebraic
automorphisms are allowed.
\begin{proof}
Suppose $P_\chi(\pi)$.  ITA5 writes
$\operatorname{algNorm}\pi$ as the ordered singleton
Langlands quotient attached to smooth characters
$\mu_0,\ldots,\mu_{n-1}$.  Their Weil restrictions
give (ITA.10), with the prescribed inertia.
Undoing the common twist, $\pi$ is the unique irreducible
quotient of normalized induction of
$\mu_i|\ |^{(n-1)/2}$.  On the upper triangular Borel,
the inducing transformation character, including the
normalizing modulus, is
\[
       \prod_{i=0}^{n-1}\mu_i(a_i)|a_i|^{n-1-i}.
                                                               \tag{ITA.13}
\]
Indeed the exponent of $a_i$ in the half-modulus is
$(n-1-2i)/2$; adding $(n-1)/2$ gives $n-1-i$.
These are integer exponents, so every value of the norm
factor in (ITA.13) is rational and fixed by $\gamma$.

Apply $\gamma$ pointwise to the functions in this induced
representation.  It is a semilinear bijection commuting
with right translation; it preserves smoothness because
stabilizers are unchanged.  Formula (ITA.13) shows that
the resulting induction is the normalized induction of
$(\gamma\mu_i)|\ |^{(n-1)/2}$.  The relation (ITA.8)
is preserved: $|\ |$ is rational-valued, so an equality
after applying $\gamma$ would give the same equality
before applying it.  Semilinear transport preserves
submodules, quotients and the unique irreducible cosocle.
Thus $\operatorname{algNorm}({}^{\gamma}\pi)$ is the
singleton quotient for $\gamma\mu_i$.

The rank-one and Langlands formulas give its parameter
$\bigoplus_i\gamma\xi_i$, with $N=0$.  Its inertia
characters are $\gamma\chi_i$.  Every value of $\chi_i$
is in $j(O)$ by ITA3, hence is fixed by $\gamma$.
This proves the forward implication in (ITA.12).
Apply the same argument to $\gamma^{-1}$ for the reverse
implication.  The cancellation in (ITA.13) is necessary:
one cannot assume that an arbitrary $\gamma$ fixes a
chosen positive square root of the residue cardinality.
\end{proof}
% END ITA7

\subsection{Descent of the detecting representation to the original DVR}

% BEGIN ITA8
\paragraph{ITA8. A stable lattice followed by restriction of scalars.}
\label{claim:ITA8}
There is a finite free, nonzero $O$-module $B$ and an action
$R:K_F\to\operatorname{Aut}_O(B)$ with open kernel such
that, with $B_j=\mathbb C\otimes_{O,j}B$,
\[
              \operatorname{Hom}_{K_F}(B_j,\pi)\ne0
                             \quad\Longleftrightarrow P_\chi(\pi)
                                                               \tag{ITA.14}
\]
for every smooth irreducible $\pi$.  The output coefficient
ring is exactly $O$.
\begin{proof}
Choose a complex basis of the finite representation $A$ of
ITA6 and its finite quotient $H$.  By ITA1 the finitely
many entries of its matrices, pulled back by $\iota^{-1}$,
belong to a finite extension $E/K$ inside the original
$\overline{\mathbb Q}_\ell$.  They define a representation
$A_E:H\to\operatorname{GL}_d(E)$ whose extension by
$\iota$ is the given $A$.  Its group identities hold over
$E$ because that field embeds in $\mathbb C$.

Let $O_E=\{x\in E:|x|_\ell\leq1\}$.  The norm-ball
argument of WRP2 proves that this is a DVR finite over
$\mathbb Z_\ell$, containing $O$.  In particular it is
finite over $O$: any finite set of $\mathbb Z_\ell$-module
generators also generates it over $O$.  Set
\[
                      \mathcal L=\sum_{h\in H}A_E(h)O_E^d
                                      \ \subset E^d.          \tag{ITA.15}
\]
The sum is finite and is preserved by every $A_E(h)$,
with equality because its inverse also belongs to $H$.
It contains $O_E^d$ and spans $E^d$ over $E$.
It is a finite torsion-free $O_E$-module, hence free of
rank $d$ over that DVR.  Regard the \emph{same} module as
an $O$-module $B=\operatorname{Res}_{O_E/O}\mathcal L$.
It is finite and torsion-free over the original DVR,
so is finite free.  Inverting $\ell$ gives $E^d$:
every element of $E$ becomes integral after multiplication
by a sufficiently high power of $\ell$.  Thus
\[
              \operatorname{rank}_O B=d[E:K]>0,
       \qquad K\otimes_O B\simeq\operatorname{Res}_{E/K}E^d.
                                                               \tag{ITA.16}
\]
The action $R$ on $B$ still factors through $H$.

Since $E/K$ is separable, scalar extension decomposes as
\[
 B_j\simeq
 \bigoplus_{\eta:E\hookrightarrow\mathbb C,\ \eta|_K=j}
                  \bigl(\mathbb C\otimes_{E,\eta}E^d\bigr).
                                                               \tag{ITA.17}
\]
For clarity, choose a primitive element of $E/K$.
Its minimal polynomial has distinct roots in $\mathbb C$;
the Chinese remainder theorem gives
$\mathbb C\otimes_K E\simeq\prod_\eta\mathbb C$
by $z\otimes x\mapsto(z\eta(x))_\eta$.  Applying this
to $E^d$ gives (ITA.17) and shows directly that each
summand's action is obtained by applying $\eta$ to the
original matrices.

Each $\eta$ extends, relative to the initial embedding
$\iota|_E$, to some $\gamma\in
\operatorname{Aut}(\mathbb C/j(K))$.  The usual extension
of an algebraic-field embedding to an algebraic closure
gives an embedding of $\mathbb C$ into itself over
$j(K)$.  It is onto: its image is algebraically closed,
and $\mathbb C$ is algebraic over $j(K)$.  Thus the
summand in (ITA.17) is ${}^{\gamma}A$.

Semilinearly transporting a map gives
\[
 \operatorname{Hom}_{K_F}({}^{\gamma}A,\pi)\ne0
 \ \Longleftrightarrow\
 \operatorname{Hom}_{K_F}(A,{}^{\gamma^{-1}}\pi)\ne0.
\]
By ITA6 and ITA7 this is equivalent to $P_\chi(\pi)$,
for \emph{every} summand of (ITA.17).  The sum is finite
and nonempty.  Therefore its Hom space is nonzero exactly
when $P_\chi(\pi)$, proving (ITA.14).
The extension $O_E$ was used only to construct the lattice;
restriction of scalars makes every final matrix an
$O$-matrix.  No assertion that $A$ itself descends
irreducibly to $O$, and no division by $|H|$ inside $O$,
is required.
\end{proof}
% END ITA8

% BEGIN ITA9
\paragraph{ITA9. The contragredient gives exactly the source's column equation.}
\label{claim:ITA9}
Choose an $O$-basis $e_0,\ldots,e_{d_B-1}$ of $B$ and let
\[
                \tau(k)=R(k^{-1})^{\mathsf T}
                       \in\operatorname{GL}_{d_B}(O).          \tag{ITA.18}
\]
This is an integral representation with open kernel.  For
every complex representation $\rho$ of $K_F$ on $M$,
the vector families satisfying the source's column equation
are naturally identified with
$\operatorname{Hom}_{K_F}(B_j,M)$.  In particular their
nonvanishing for a smooth irreducible $\pi$ is equivalent
to $U_\chi(\pi)$ of ITA5.
\begin{proof}
The module underlying (ITA.18) is $B^*=\operatorname{Hom}_O(B,O)$,
with the dual basis.  It is preserved because $R(k)$ and
$R(k^{-1})$ preserve $B$.  Matrix multiplication gives a
homomorphism, and the kernel equals that of $R$.
To a family $F_a\in M$ associate the unique linear map
$f:B_j\to M$ with $f(e_a)=F_a$.  For every $k,a$,
\[
 \sum_b j(\tau(k)_{ab})\rho(k)F_b
       =\rho(k)f\bigl(R(k^{-1})e_a\bigr).                   \tag{ITA.19}
\]
This equals $F_a$ for all $a$ exactly when
$fR(k)=\rho(k)f$.  Thus the family's equations are
precisely equivariance, and the family is nonzero exactly
when $f$ is nonzero.  Equivalently the family represents
an invariant tensor in $B_j^*\otimes_{\mathbb C}M$.
Apply ITA8 and ITA5 for the last assertion.  This explicitly
checks the inverse and transpose; using $R$ itself in
the column equation would detect its dual instead.
\end{proof}
% END ITA9

\subsection{The product of local types for arbitrary isotypic modules}

% BEGIN ITA10
\paragraph{ITA10. What the literal \texttt{IsoSum} hypothesis supplies.}
\label{claim:ITA10}
Let $G$ be a group, $\pi$ an irreducible complex representation,
and $N$ a representation generated by images of equivariant
maps $\pi\to N$.  Then $N$ is a direct sum of copies of
$\pi$.  Every subrepresentation of such an $N$ is again
a sum of copies of $\pi$.  If $N\ne0$ and $C$ is a
finite-dimensional representation of a subgroup $K\subset G$,
then
\[
                 \operatorname{Hom}_K(C,N)\ne0
                   \quad\Longleftrightarrow\quad
                 \operatorname{Hom}_K(C,\pi)\ne0.            \tag{ITA.20}
\]
\begin{proof}
A nonzero map from $\pi$ is injective, and its image is
simple.  Choose a maximal family of these simple
submodules whose sum is direct.  If its sum $N_0$ is not
all of $N$, some generating simple submodule $P$ is not
contained in $N_0$.  Its intersection with $N_0$ is a
submodule of the simple $P$, hence is zero.  Adding $P$
contradicts maximality.  This proves the first statement.

For completeness the same argument is the usual
semisimple-module argument for submodules.  A sum of
simple modules is semisimple: given a submodule $Y$, take
a maximal direct sum of simple summands disjoint from
$Y$; its sum with $Y$ must be all of $N$, because a
simple generating summand not contained in that sum has
zero intersection with it and could be added.  Hence
every submodule has a complement.  Project the original
simple generating summands to $Y$ along that complement.
Their nonzero images are simple and isomorphic to
$\pi$, and generate $Y$.  This proves the asserted
isotypy of submodules, without a finite length assumption.

Finally a linear map from finite-dimensional $C$ into a
direct sum has image in a finite subsum: apply the
definition of direct sum to the images of a basis.
If a nonzero equivariant map $C\to N$ exists, one
coordinate projection to a copy of $\pi$ is nonzero
and equivariant.  Conversely, since $N\ne0$, there is
a copy of $\pi$ in $N$ into which any nonzero map
$C\to\pi$ can be composed.  This proves (ITA.20).
\end{proof}
% END ITA10

% BEGIN ITA11
\paragraph{ITA11. Simultaneous detection in the original product representation.}
\label{claim:ITA11}
For each $u\in S$ take the $O$-module $B_u$ of ITA8,
its action $R_u$, and its dual action $\tau_u$ from ITA9.
Put
\[
 B_S=\bigotimes_{u\in S,O}B_u,\qquad
 R_S(k)=\bigotimes_{u\in S}R_u(k_u),\qquad
 \tau_0=R_S^*,\qquad m=\prod_{u\in S}\operatorname{rank}_O B_u.
                                                               \tag{ITA.21}
\]
Choose the tensor basis and any fixed identification of its
finite index set with $\{0,\ldots,m-1\}$.  For every
$M,\rho,(\pi_u)$ satisfying the exact hypotheses before
(ITA.6), its left side holds if and only if
$P_{\chi_u}(\pi_u)$ holds for every $u$.
\begin{proof}
By ITA9 the left side is the nonvanishing of
\[
 \bigl((B_S)_j^*\otimes M\bigr)^{K_S}
                \simeq\operatorname{Hom}_{K_S}((B_S)_j,M).
                                                               \tag{ITA.22}
\]
Enumerate the finite set $S$.  Starting with $N_0=M$,
apply successively the functors
\[
                    N\longmapsto
                    \operatorname{Hom}_{K_u}((B_u)_j,N)
                      \simeq((B_u)_j^*\otimes N)^{K_u}.
                                                               \tag{ITA.23}
\]
The remaining groups act on the $N$ factor.  These actions
are well defined and commute with taking $K_u$-invariants,
since the original factors of $G_S$ commute.  For each
unprocessed factor $G_v$, the previous module is
$\pi_v$-isotypic: this is true at the start by the
literal \texttt{IsoSum} hypothesis and ITA10.  The next
module is a $G_v$-submodule of the finite direct sum
$(B_u)_j^*\otimes N$, and so is again $\pi_v$-isotypic
by ITA10.

Whenever the previous module $N$ is nonzero, (ITA.20)
and (ITA.14) say that (ITA.23) is nonzero exactly when
$P_{\chi_u}(\pi_u)$ holds.  A zero module stays zero
under all subsequent steps.  Induction therefore gives
nonzero output exactly when every one of these local
conditions holds.  Successive finite tensor--Hom
identifications identify that output with (ITA.22):
equivariance in each coordinate is precisely equivariance
under their finite product.  This proves the assertion.
No decomposition of $M$ as an external tensor product,
no scalar-endomorphism version of Schur's lemma, and no
admissibility of $M$ have been used.

If $S$ is empty, (ITA.21) means $B_S=O$, $m=1$ and
the trivial action.  Then (ITA.22) is $M\ne0$, and
the right side is the empty universal condition, so the
same conclusion holds in this case as well.
\end{proof}
% END ITA11

% BEGIN ITA12
\paragraph{ITA12. The complete original integral \texttt{TypeForAlt}.}
\label{claim:ITA12}
The homomorphism $\tau_0$ in (ITA.21), over the unchanged
ring $O=c.O$, satisfies the entire original
$\texttt{TypeForAlt}(r,\iota,\tau_0)$.
\begin{proof}
First the intersection of the coordinate inverse images
of $\ker R_u$ is open in $K_S$.  Every element of that
intersection acts trivially on the tensor module in
(ITA.21); hence it belongs to $\ker\tau_0$.
A subgroup containing an open subgroup is open, since
it is a union of its open cosets.  Its image under the
injective scalar map $j$ has exactly the same kernel:
a matrix becomes the identity over $\mathbb C$ if and
only if it was the identity over $O$.  Thus the
complexified kernel in the source's \texttt{IsTypeAlt}
field is open.

Fix an arbitrary family $t'$ satisfying the source's
$r.\texttt{IsTeich}$, not necessarily the family used
when choosing the local types.  ITA2 gives equality
$t'_i=t_i$ as homomorphisms for every $i$.  Consequently
the local characters (ITA.5) coincide as actual maps
on the original inertia subgroups.  For every allowed
$M,\rho,(\pi_u)$, ITA11 proves the left side of
(ITA.6) equivalent to all the local conditions
$P_{\chi_u}(\pi_u)$.  ITA5 turns each of them into
the universally quantified \texttt{IsRecT} condition.
The coordinate formula in ITA9 is exactly the sum in
the source definition.  This gives every field of
\texttt{IsTypeAlt} for the fixed $t'$, and $t'$ was
arbitrary.  Therefore it proves \texttt{TypeForAlt}.

The choice order is the required one.  The data
$s,n,V,Q,c,r,\iota$ are fixed first; then the unique
Teichm\"uller lifts, finitely many local types and
lattices, their $O$-bases, and their tensor product
are chosen once.  The resulting $m,\tau_0$ precede
every later lift family, local correspondence, module,
and representation in (ITA.6).  The only coefficient
extension was internal to (ITA.15), followed by
restriction of scalars.  All output matrices lie in
the original $O$ and are complexified through the
specified $\iota$.

This proves the existential endpoint of
\path{TTypeEx.exists_integral_type_for_residual_rep}
under its original foundational inputs.  Its rank is
every original $n\geq1$.  The argument includes the
allowed original coefficient prime $\ell=3$, every
local residue prime, repeated inertia characters,
and the empty type-place set.  It imposes no
genericity, regularity, tame-ramification, large-prime
or prime-to-rank hypothesis.
\end{proof}
% END ITA12

% BEGIN ITA13
\paragraph{ITA13. The precise input supplied to the existing Hecke-data assembly.}
\label{claim:ITA13}
For the same original $V,Q,c,r,\iota$ and every actual
universal character ring $\Lambda$, ITA12 supplies the
formerly missing type input in OTW13.  Applying OTW13
with this $m,\tau_0$ gives the original
$\texttt{UnitaryHeckeData}(Q,\iota,\Lambda)$, and hence
its \texttt{HdEx} existential conclusion.
\begin{proof}
The integral representation and all its maps have the
types required by OTW13; ITA12 gives precisely its
\texttt{type} hypothesis.  OTW12 in that assembly
provides the one character $\theta$ together with
\texttt{theta} and \texttt{isWt}, using this same
representation.  CAP17 provides the residual
Hecke-eigenvalue function $a$ and its \texttt{res}
field on the same original operators.  Thus
\[
        (m,\tau_0,\texttt{type},\theta,
                   \texttt{theta},\texttt{isWt},a,\texttt{res})
\]
has exactly the eight fields of
\texttt{UnitaryHeckeData}.  This is direct substitution
into the already reconstructed OTW13 proof, with no
new coefficient extension or later alteration of
$\tau_0$, $r$ or $\Lambda$.  The conclusion is the
Hecke-data supplier only; the later common-ideal
and automorphy conclusions require their own
separate arguments.
\end{proof}
% END ITA13

\subsection{Proof boundary and exact source dictionary}

The following locators are in the pinned source commit
\begin{center}\texttt{7cbf7a0c32fca1badf05931879a37c3f8a2b03b4}.\end{center}
The companion metadata records their full file and range hashes.
\begin{center}\small
\begin{tabular}{p{0.26\linewidth}p{0.38\linewidth}p{0.26\linewidth}}
\hline
Original object or endpoint & Source file and declaration & Written proof \\
\hline
Actual coefficient subring & \texttt{DirichletPrimeToLPart}\allowbreak\texttt{MulLPowerPart.Coef} & ITA1, ITA8 \\
Every lift family & \path{ArithmeticPointClassicality.Resid.IsTeich} & ITA2--ITA3 \\
Actual split type places & \path{ContinuousResP.HidaSetting}, \texttt{Lev.pr} & ITA3--ITA4 \\
Zero-monodromy diagonal inertia & \texttt{Distinct.IsTypeWD} & ITA5 \\
Weaker local correspondences & \texttt{IsRecTOfIsRec.LawsT}, \texttt{IsRecT} & ITA4--ITA5 \\
Conjugation and integral descent & \texttt{TypeConj}, \texttt{TypeLoc}, \texttt{PTypeEx.te\_loc} & ITA7--ITA9 \\
Arbitrary product modules & \texttt{CoefficientRep}\allowbreak\texttt{TrivialOnUnramified.IsoSum} & ITA10--ITA11 \\
Full original type & \path{UnitaryHeckeData.TypeForAlt}, \texttt{TTypeEx} & ITA12 \\
Original Hecke-data record & \texttt{UnitaryHeckeData}, \texttt{HdEx} & ITA13, OTW13 \\
\hline
\end{tabular}
\end{center}

The new proof simplifies the source's \texttt{NonzeroTypes}
and generic integral-annihilator descent at the existential
endpoint actually required.  It does not claim that the chosen
lattice basis equals the source's auxiliary Hilbert-choice
basis, and no such equality is required by that endpoint.
The local type theorem is an explicit matched classical
input, not a new proof of all of Bushnell--Kutzko or
Schneider--Zink theory.  Likewise \texttt{RecTExists}
is retained at its original foundation status.  The
mathematical progress here is the full coefficient,
normalization, universal-quantifier, duality and product
argument yielding ITA12 and its exact OTW13 consequence.
